\cotefile{157-2}
\scanpage{115}\mspage{60}
sur \unsure{$p$|$s$} (\unsure{au niveau local|au-dessus local}), donc le chgt de base ${}_{s'\backslash}S'\unsure{\longrightarrow|\text{---}}S$ s'identifie à ${}_{s\backslash}S\unsure{\text{---}|\longrightarrow}S$, et l'image inverse est bien isom à ${}_{s\backslash}X$, OK.

c) Que de \unsure{$\overline{\mathrm m}$}
\[
 \Phi\,\mathrm{id}\longrightarrow\mathrm{id}
\]
est non seulement dans $W_S^{\deleted{\unsure{s}}}$ (grossier), mais que c'est une $W$-équivalence \deleted{\unsure{fibr}} par fibres. Donc on part de $X$ \unsure{fibrée|fibré} sur $S$, on regarde
\[
 \Phi(X)\longrightarrow X
\]
pour voir que c'est une $W$-équiv. par fibres, il suffit de voir qu'il en est ainsi de $X\longrightarrow\Phi(\unsure{x|X})$, \deleted{\unsure{Or} \illegible{2}} \inserted{\unsure{en|car} \unsure{pos|pour}}
\[
 X_s\longrightarrow{}_{s\backslash}X
\]
est une $W$-équiv. pour \unsure{tt|$\forall$} $s\in S$. Mais ceci est déjà vrai dès que $X$ est lisse sur $S$.

Comme le cas de la \unsure{$W$|bi}-équivalence fine gauche se ramène au cas actuel par $S\mapsto S^\circ$, on a donc prouvé tout ce qui était attendu, à la seule exception de la pleine fidélité \unsure{en|des} foncteurs $i,j$ de la page 51.

\medskip
Mais dans le résumé prospectif de la situation pour les \unsure{localiseurs}, p.\,\unsure{49|59}-51, j'ai oublié d'autres foncteurs que $i,j$, de nature si possible encore plus triviale. Ce sont
% archive sheet number: 114
\scanpage{116}\mspage{61}
ceux qu'on peut \unsure{décrire|dériver} à partir du foncteur identique, sur une des catégories du diagramme \deleted{\unsure{que}} où aient cours deux, \inserted{voire trois,} notions d'\inserted{$W$-}équivalences différentes. Ce sont, je le rappelle

$1^\circ$) Les \deleted{\unsure{c}} quatre catégories
\[
 \text{\uline{Fibsc}}\,S\unsure{\hookrightarrow|\longrightarrow}\text{\uline{Fib}}\,S\longrightarrow\text{\uline{Liss}}\,S\unsure{\longrightarrow|\text{---}}\mathrm{Cat}_{/S},\quad W_S^d\subset W_S
 \tag*{$(*)$}
\]
pour \unsure{la|les} \deleted{\unsure{éq}} $W$-\unsure{équivalence|équivalences} fine droite d'une part, grossière de l'autre, et dualement
\[
 \text{\uline{Cofibsc}}\,S\longrightarrow\text{\uline{Cofib}}\,\deleted{\unsure{sc}}S\longrightarrow\text{\uline{Prop}}\,S\longrightarrow\mathrm{Cat}_{/S},\quad W_S^g\subset W_S
 \tag*{$(*')$}
\]
pour la $W$-équiv. fine gauche d'une part, la grossière de l'autre\textup{;} et \deleted{\unsure{du même}} \inserted{de plus}

$2^\circ$) $\mathrm{Cat}_{/S}$, qui était \unsure{incluse|inclus} dans les deux suites précédentes, et qui est la seule cat. du diagramme donnant lieu à trois notions de $W$-équivalences différentes.
\diagram{\begin{tikzcd}[ampersand replacement=\&,column sep=small,row sep=0pt]
 W_S^d\arrow[dr,phantom,"{\subset}" sloped] \& {}\\
 {} \& W_S\,.\\
 W_S^g\arrow[ur,phantom,"{\subset}" sloped] \& {}
\end{tikzcd}}

Dans chacun des quatre cas \deleted{$\unsure{C,\Sigma,\Sigma}$} ($C$, munie des deux localiseurs $\Sigma\subset\Sigma'$) de la suite $(*)$, \deleted{pour voir la nature} \deleted{\unsure{du}} le foncteur
\[
 C\Sigma^{-1}\longrightarrow C\Sigma'^{-1}
\]
est essent. ``le même'', à cause des th. de comparaison. Il suffit de regarder ce foncteur dans un seul cas, p. ex. un des deux extrêmes $\text{\uline{Fibsc}}\,S$ ou $\mathrm{Cat}_{/S}$. Même remarque pour la suite $(*')$. Donc finalement il reste deux foncteurs remarquables
% archive sheet number: 115
\scanpage{117}\mspage{62}
à examiner, qui n'ont pas l'air du tout de provenir des foncteurs du diagramme p.~8 ou de quelque composé (ça ne donne que des équivalences canoniques ou bien les foncteurs $i,j$ de page 51\unsure{)|,}, quand on va des cat. internes \unsure{au|en} diagramme vers les catégories \unsure{bordantes}). Le plus joli est sans doute de les regarder à partir de $\mathrm{Cat}_{/S}$, puisque c'est là seulement qu'on les a tous les deux à la fois
\diagram{\begin{tikzcd}[ampersand replacement=\&,column sep=small,row sep=small]
 \mathrm{Cat}_{/S}(W_S^d)^{-1}\approx\mathrm{HOT}_W(S)\arrow[dr,"\unsure{\rho|p}"] \& {}\\
 {} \& \mathrm{Cat}_{/S}W_S^{-1}\approx\mathrm{HOT}_W^{\unsure{\mathrm{lc}|\ell\mathrm{c}}}(S)\\
 \mathrm{Cat}_{/S}(W_S^g)^{-1}\approx\mathrm{HOT}_W^\circ(S)\arrow[ur,"\sigma"] \& {}
\end{tikzcd}}

Par contre, les foncteurs (\unsure{qu'on a envie|que j'ai envie} d'appeler ``foncteurs d'\unsure{induction|inclusion}'') $i,j$ de la page 51, en sens inverse de $\rho,\sigma$, sont ceux qui proviennent, par localisation $W_S^d$ ou $W_S^g$ selon les cas, d'un des foncteurs du diagramme page 8 (ou d'un composé de tels foncteurs), allant d'une des dix catégories internes (où toutes les $W$-équivalences envisagées coïncident) vers l'une des sept catégories \unsure{bordantes}, sauf quand cette dernière est $\mathrm{Cat}_{/S}$, et détermine déjà avec quelle $W$-équivalence, $W_S^d$ ou $W_S^g$, on va travailler. La détermi%
% archive sheet number: 116
\scanpage{118}\mspage{63}%
nation peut-être la plus jolie est par
\[
 \text{\uline{Parf}}\,S\longrightarrow\mathrm{Cat}_{/S}
 \tag*{$(*)$}
\]
qui donne le foncteur $i$ quand on travaille avec $W_S^d$, le foncteur $j$ quand on travaille avec $W_S^g$. On voit \inserted{\unsure{alors|donc}} tout de suite que si on compose avec le foncteur pertinent, soit $\unsure{\rho|p}$ soit $\sigma$, défini par l'inclusion $W_S^{\unsure{d|g}}\unsure{\in|\subset}W_S$ resp $W_S^g\subset W_S$ (comme parties de $\mathrm{Fl}(\mathrm{Cat}_{/S})$), on trouve le foncteur \deleted{canonique}
\[
 (\text{\uline{Parf}}\,S)W_S^{-1}\overset{\approx}{\unsure{\longrightarrow|\text{---}}}(\mathrm{Cat}_{/S})W_S^{-1}
 \tag*{$\unsure{(**)|(*)'|(*)}$}
\]
qui est l'\uline{équivalence canonique} entre ces deux catégories. Ainsi le \deleted{\unsure{des}} foncteur correspondant à $i,j$ page 51, exprimé en termes de $\mathrm{Cat}_{/S}$ comme foncteur en sens inverse de $\rho,\sigma$ (p.~62), est-il le composé de l'équiv. quasi-inverse de \unsure{$(**)$}, laquelle s'explicite par le foncteur \unsure{$\mathcal C$|$C$}, et du foncteur sur les \unsure{localisés|localis.} $W_S^d$ ou $W_S^g$ déduit de $(*)$, i.e.
\diagram{\scriptsize\begin{tikzcd}[ampersand replacement=\&,column sep=small,row sep=large]
 {} \& {} \& \mathrm{HOT}_W(S)\approx(\mathrm{Cat}_{/S})(W_S^d)^{-1}\\
 \mathrm{HOT}_W^{\unsure{\mathrm{lc}|\ell\mathrm{c}}}(S\unsure{.|{}})\approx(\mathrm{Cat}_{/S})W_S^{-1}
 \arrow[urr,"i"]\arrow[rrd,"j"']\arrow[r,"{\deleted{\unsure{\approx}}\mathcal C}"]
 \&(\text{\uline{Parf}}\,S)W_S^{-1}
 \arrow[ur,"{\text{car }W_S=W_S^d\text{ sur \uline{Parf}}}"']
 \arrow[dr,"{\text{car }W_S=W_S^g\text{ sur \uline{Parf}}}"]
 \&{}\\
 {} \& {} \& \mathrm{HOT}_W^\circ(S)\approx(\mathrm{Cat}_{/S})(W_S^g)^{-1}\,.
\end{tikzcd}}
% archive sheet number: 117
\scanpage{119}\mspage{64}
On voit par là qu'on a des isom. canoniques
\[
 \unsure{\rho|p}\,i\simeq\mathrm{id},\qquad \sigma j\simeq\mathrm{id}
\]
dans la catégorie des endofoncteurs de $(\mathrm{Cat}_{/S})W_S^{-1}$ $(\approx\mathrm{HOT}^{\unsure{\mathrm{lc}|\ell\mathrm{c}}}_W(S))$, \deleted{\unsure{puisque}} isomorphismes qui s'explicitent ici par \deleted{\unsure{l'isom} \unsure{homotopique|homotopie}} \inserted{le fait} que $\rho i,\sigma j$ sont déduits du foncteur $\unsure{C|\mathcal C}$ par passage \unsure{aux localisés|au localisé} \inserted{$W_S^{-1}$} (donc en fait $\rho i=\sigma j$ !), et qu'on a une \deleted{\unsure{isom}} flèche fonctorielle en $X$
\[
 X\xrightarrow{\unsure{i_x|i_X}}C(X),
\]
qui est dans $W$ pour tt $X$. Ainsi, $\rho$ et $\sigma$ sont \inserted{à isom près} des \uline{rétractions} des \deleted{\unsure{C}} catégories $\mathrm{HOT}_W(S)$, $\mathrm{HOT}_W^\circ(S)$ sur la catégorie $\mathrm{HOT}_W^{\mathrm{lc}}(S)$. \deleted{Celle-ci} De \deleted{est \unsure{équivale}} plus, ces rétractions sont \deleted{\unsure{des}} \deleted{\unsure{c. à dire|c'est à dire}} ce qu'on appelle des foncteurs de localisation (comme on voit sur le diagramme de la page 62). (On dit que $C\xrightarrow{F}C'$ est un \uline{foncteur de localisation} si, désignant par $\Sigma\subset\mathrm{Fl}\,C$ l'ens des flèches de $C$ transformées en des iso par $F$, le foncteur
\[
 C\Sigma^{-1}\xrightarrow{\overline F}C'
\]
qui factorise $F$ est une équivalence de catégories.)
% archive sheet number: 118
\scanpage{120}\mspage{65}
On sait d'ailleurs que si un foncteur \inserted{$f$} a un adjoint $g$ (à gauche ou à droite), alors $f$ est un foncteur localisation \unsure{sss|ssi} l'adjoint $g$ est pl. fidèle. Pour prouver que $i,j$ sont pl. fidèles, il serait donc tentant de prouver qu'ils sont adjoints des foncteurs $\unsure{\rho|p},\sigma$.

Pour \deleted{\unsure{voir}} y \unsure{venir à bout|voir un peu plus clair}, je vais poser
\[
 \left\{\begin{aligned}
 \deleted{\unsure{\mathcal A}}\unsure{A|\mathcal A}&=\mathrm{Cat}_{/S},\quad
 \deleted{\unsure{W_0}}\Sigma_0=W_S^d,\quad\Sigma=W_S\qquad(\subset\mathrm{Fl}(A))\\
 C:\ A&\unsure{\longrightarrow|\text{---}}A\quad\text{le foncteur familier, qui a la vertu}\\
 &C(\Sigma)\subset\Sigma_0
 \end{aligned}\right.
\]
On a donc la situation
\diagram{\begin{tikzcd}[ampersand replacement=\&,column sep=large]
 {(\unsure{\alpha|d})}\& (A,\Sigma_0)\arrow[r,shift left,"\rho_0=\mathrm{id}_A"]\&
 (A,\Sigma),\arrow[l,shift left,"C"]
\end{tikzcd}}
induisant par passage aux localisés
\diagram{\begin{tikzcd}[ampersand replacement=\&,column sep=large]
 A\Sigma_0^{-1}\arrow[r,shift left,"\rho=\overline{\rho_0}"]\&
 A\Sigma^{-1},\arrow[l,shift left,"\overline C"]
\end{tikzcd}}
on veut \deleted{\unsure{dégager}} vérifier que $\rho$ et $C$ sont adjoints l'un de l'autre, $C$ \uuline{sûrement à droite}, ce qui signifie que l'on doit avoir des hom fonctoriels
% slanted
\margin{mais si $C$ est à \uline{droite} les flèches ici doivent aller en sens inverse et c'est bien à quoi je dois \unsure{en venir|me résoudre} \unsure{on dirait|on dit} p.~67}
\diagram{\begin{tikzcd}[ampersand replacement=\&,row sep=small,column sep=small]
 {}\arrow[d,phantom,"{\Bigg\{}" description] \&
 \mathrm{id}_{A\Sigma^{-1}}\arrow[r,shift left]\&
 \rho\overline C\arrow[l,shift left,very thick]\& (= \overline{\rho_0 C})\\
 {}\&\mathrm{id}_{A\Sigma_0^{-1}}\arrow[r,shift right,very thick]\&
 \overline C\rho\arrow[l,shift right]\&\unsure{\neq|{}}(=\overline{C\rho_0})
\end{tikzcd}}
satisfaisant deux relations de compatibilité sur lesquelles on reviendra. On peut espérer trouver ces \deleted{\unsure{foncteurs}} hom fonctoriels à partir de
\[
 \left\{\begin{aligned}
 \mathrm{id}_A&\xrightarrow{i}\rho_0 C=C\\
 \mathrm{id}_A&\xleftarrow{\unsure{p|{}}} C\rho_0=C
 \end{aligned}\right.
\]
% archive sheet number: 119
\scanpage{121}\mspage{66}
et dans notre cas on dispose bel et bien de
\[
 i:\mathrm{id}_{\unsure{A|\mathcal A}}\longrightarrow C\qquad\text{i.e. de }i_X:X\longrightarrow C(\unsure{x|X})
\]
fonctoriel en $X\in\mathrm{Ob}\,A$ (et de plus $i_X$ est dans $\Sigma$). Il n'y a pas cependant de flèche naturelle en sens inverse $p_X$
\[
 C(X)\xrightarrow{?}X.
\]
J'avais songé à la projection $\pi_X:C(X)\longrightarrow X$, déduite de la construction $C(X)\simeq X\times_S(\text{\uline{Ch}}_{\unsure{\infty|w}}(S),\sigma)$, mais ce n'est malheureusement pas un $S$-morphisme. Comme on cherche une flèche de $A\Sigma_0^{-1}$, on pourrait se contenter d'avoir une flèche fonctorielle en sens inverse, du moment qu'elle soit dans $\deleted{\unsure{W}}\Sigma_0$, donc puisse s'inverser quand on travaille dans $A\Sigma_0^{-1}$. Le candidat évident est $i_X:X\unsure{\longrightarrow|\text{---}}C(X)$ dont on vient de \unsure{parler}. L'ennui, c'est que $i_X$ est seulement dans $\Sigma$, pas dans $\Sigma_0$.\textsuperscript{*}
% slanted
\margin{s'il n'en était pas ainsi, le type d'homotopie relatif sur $S$ défini par $X/S$ serait localement constant (puisqu'il en est ainsi \unsure{pour celui} défini par $C(X)$), ce qui \unsure{n'est pas} tjrs le cas.}
Pourtant, quand $X$ est lui-même de la forme $C(Y)$ (plus généralement, s'il est un \unsure{Hot fibré|$\mathrm{Hot}_W$ fibré} sur $S$) alors \deleted{\unsure{$X$}} le fait que $X\unsure{\text{---}|\longrightarrow}C(X)$ soit dans $\Sigma$ implique déjà qu'il soit dans $\Sigma_0$. Notons donc ceci
\[
 \left\{\begin{aligned}
 i:\mathrm{id}_A\longrightarrow C\quad\text{tel que}\quad&
 1^\circ)\quad i_X:X\longrightarrow C(X)\quad\text{\uline{dans $\Sigma$}}\\
 &\hspace{2em}\text{pour \unsure{tt|$\forall$} }X\in\mathrm{Ob}\,A\\
 &2^\circ)\quad i_X\ \text{\deleted{dans}}\ \Sigma_0\quad\text{si }X\text{ de la}\\
 &\hspace{2em}\text{forme }C(Y),
 \end{aligned}\right.
\]
\deleted{On va poser alors
\[
 \deleted{\unsure{C}}\,A_0=C(A)\subset A_0,
\]
on va considérer les foncteurs
\diagram{\begin{tikzcd}[ampersand replacement=\&]
 A_0\arrow[r,hook,shift left,"\rho_0"]\& A\arrow[l,shift left,"C"]
\end{tikzcd}}
$\rho_0$, l'inclusion}
% archive sheet number: 120
\scanpage{122}\mspage{67}
L'idée me vient, après des essais \unsure{embrouillants|embrouillés}, que $\overline C$ pourrait être l'adjoint à \deleted{gauche}\inserted{droite}, et non à \deleted{droite}\inserted{gauche}, i.e. que les flèches cherchées pourraient aller en sens inverse
\[
 \begin{aligned}
 \mathrm{id}_{\unsure{A|\mathcal A}\Sigma^{-1}}&\xleftarrow{p}\unsure{\rho|p}\,\overline C\quad(=\overline{\rho_0 C}\deleted{\unsure{*}})\\
 \mathrm{id}_{A\Sigma_0^{-1}}&\xrightarrow{i}\overline C\rho\quad(=\overline{C\rho_0})
 \end{aligned}
\]
(j'inverse du coup les notations pour $p$ et $i$), donc on est amené si possible à chercher des flèches
\[
 \begin{aligned}
 \mathrm{id}_A&\mathrel{\deleted{\unsure{\rightarrow}}\!\!\xleftarrow{p}}\rho_0 C=C\\
 \mathrm{id}_A&\xrightarrow{i}C\rho_0=C,
 \end{aligned}
\]
donc des flèches fonctorielles en $X$
\[
 X\xrightarrow{i_X}C(X)\xrightarrow{p_X}X.
\]
On tient bel et bien un \deleted{$\unsure{p_X}$} $i_X$ fonctoriel --- tjrs le même ! Pour $p_X$, on n'en a pas plus qu'avant, mais se rappelant qu'il suffit de le trouver dans $A\Sigma^{-1}$ (et non plus $A\Sigma_0^{-1}$ !), il suffit d'inverser la flèche $i_X$, qui est en effet un \uline{iso} dans $A\Sigma^{-1}$ (mais non dans $A\Sigma_0^{-1}$). On posera donc
\[
 p_X=i_X^{-1}:C(X)\unsure{\text{---}|\longrightarrow}X\qquad\text{dans }A\Sigma^{-1}.
\]
Il faut prouver deux compatibilités seulement, pour les foncteurs $\overline C$ et $\rho$ entre $A\Sigma_0^{-1}$ et $A\Sigma^{-1}$ (et j'\deleted{\unsure{é}}écrirai $C$ au lieu de $\overline C$)\,: les deux composés
\diagram{\scriptsize\begin{tikzcd}[ampersand replacement=\&,column sep=small,row sep=large]
 {}\arrow[d,phantom,"{\Bigg\{}" description]\& \rho\arrow[r,"{\rho*i}"]\arrow[rr,bend left=28,dashed,"\mathrm{id}"]\&
 \rho C\rho\arrow[r,"{p*\rho}"]\&\rho\&\text{(relation \deleted{\unsure{pour les flèches|pour foncteurs}} à valeurs dans } A\Sigma^{-1}\text{)}\\
 {}\& C\arrow[r,"{i*C}"]\arrow[rr,bend right=23,dashed,"\mathrm{id}"']\&
 C\rho C\arrow[r,"{C*p}"]\& C\&\text{(relation pour foncteurs à valeurs dans } A\Sigma_0^{-1}\text{)}
\end{tikzcd}}
doivent être l'identité $\mathrm{id}_{\rho}$ resp $\mathrm{id}_C$.
% archive sheet number: 121
\scanpage{123}\mspage{68}
La première relation se vérifie tautologiquement, elle revient à dire que dans $\unsure{A|\mathcal A}\Sigma^{-1}$, on a un composé identité de
\[
 X\overset{i_X}{\unsure{\text{---}|\longrightarrow}}C(\unsure{x|X})\overset{i_X^{-1}}{\text{---}}X\qquad !
\]
La seconde signifie que pour \unsure{tt|$\forall$} $x$, on a un composé identité dans $A\Sigma_0^{-1}$
\[
 C(X)\xrightarrow{i_{C(X)}}C(C(X))\xrightarrow{C(p_X)}C(X)
\]
Or $p_X:C(X)\text{---}X$ avait été défini dans $A\Sigma^{-1}$ \deleted{\unsure{que}} comme inverse $i_X^{-1}$ de $i_X$, d'autre part le \deleted{$\unsure{C}$} \inserted{foncteur $C$ sur $A$ \unsure{définit|défini} par \unsure{hyp}} $A\Sigma^{-1}\text{---}A\Sigma_0^{-1}$, plus précisément $C(\Sigma)\subset\Sigma_0$, ainsi $C(i_X)$ est inversible dans \deleted{$\unsure{C\,C}$} $A\Sigma_0^{-1}$, et $C(p_X)=C(i_X)^{-1}$. Donc la compatibilité revient à la relation
\[
 C(i_X)^{-1}\deleted{\unsure{C}}\,i_{C(X)}=\mathrm{id}\qquad\text{\uline{dans $A\Sigma_0^{-1}$}}
\]
i.e.
\[
 \boxed{i_{C(X)}=C(i_X)}:\quad C(X)\longrightarrow C^2(X).
\]
En d'autres termes, il faut prouver le

\noindent\begin{tabular}{@{}|p{.94\linewidth}@{}}
\uline{Lemme} \quad Soit $X$ dans $\mathrm{Cat}_{/S}$, considérons les deux flèches
\[
 i_{C(X)},\ C(i_X):\quad C(X)\unsure{\rightrightarrows|\substack{\text{---}\\\text{---}}}C(C(X)).
\]
Ces flèches sont égales, sinon \deleted{\unsure{du}}au sens absolu, du moins dans $(\mathrm{Cat}_{/S}\unsure{]|)}(W_S^d)^{-1}$.
\end{tabular}

\deleted{\unsure{Défi} \unsure{de}}
\[
 \begin{aligned}
 C(X)&=\{x,c\mid x\in\mathrm{Ob}\,X,\ c:f_X(x)\underset{\substack{\text{chemin}\\\text{dans }S}}{\longrightarrow}t\}\\
 &\hspace{4em}(f_X:X\unsure{\longrightarrow|\text{---}} S\ \text{la proj. str.}\\
 C^2(X)&=\{x,c,c'\mid x\in\mathrm{Ob}\,X,\ \deleted{\unsure{c:}}f_X(x)\xrightarrow{c}t\overset{c'}{\longrightarrow}t'\}\\
 &\hspace{4em} c,c'\ \text{chemins dans }S,
 \end{aligned}
\]
les projections dans $S$ sont données par
\[
 \begin{aligned}
 (x,c)&\longmapsto t=\mathrm{but}\,c\\
 (x,c,c')&\longmapsto t'=\mathrm{but}\,c'.
 \end{aligned}
\]
% archive sheet number: 122
\scanpage{124}\mspage{69}
On trouve \deleted{\unsure{pour} \unsure{$i$}}
\diagram{\begin{tikzcd}[ampersand replacement=\&,row sep=small,column sep=small]
 {} \&x\arrow[d,no head]\&{}\& {}\&\text{chemin vide}\\
 i_{C(X)}(x,c)=(x,c,1_t):\&f_X(x)\arrow[r,"c"]\&t\arrow[r,"{1_t}"{name=one}]\&t\arrow[from=one,to=1-5]\&{}
\end{tikzcd}}
\diagram{\begin{tikzcd}[ampersand replacement=\&,row sep=small,column sep=small]
 {}\&x\arrow[d]\&{}\&{}\\
 C(i_X)(x,c)=(x,1_{f_X(x)},c):\&f_X(x)\arrow[r,no head,"{1_{\unsure{f(x)|f_X(x)}}}"{name=empty}]\&f_X(x)\arrow[r,"c"]\&t\\
 {}\&{}\&\text{chemin vide}\arrow[from=empty,to=3-3]\&{}
\end{tikzcd}}

Mais on a le choix entre deux théories des chemins, pour construire $C(X)$\,: soit avec les $\text{\uline{ch}}_{\infty}(S)$, soit avec les $\text{\uline{Ch}}(S)$, ``Ça donne les mêmes résultats'', mais suivant les situations l'un est plus commode que l'autre. Je vais travailler avec \unsure{``les''|les} $\text{\uline{Ch}}(S)$, auquel cas les $1_s$ sont \deleted{des} unités à gauche et à droite, \deleted{\unsure{on}} et on a une loi de composition pour les chemins. Cette loi nous permet de définir
\[
 \begin{aligned}
 C^2(X)&\xrightarrow{\unsure{r|\gamma}}C(X)\\
 (x,c,c')&\longmapsto(x,c'\circ c),
 \end{aligned}
\]
et les composés de $i_{C(X)}$ et de $\unsure{C_{i(X)}|C(i_X)}$ avec cette flèche sont $\mathrm{id}_{C(X)}$ dans les deux cas
\diagram{\begin{tikzcd}[ampersand replacement=\&]
 C(X)\arrow[r,shift left]\arrow[r,shift right]\arrow[rr,bend right=27,"\mathrm{id}"']\&
 C^2(X)\arrow[r,"r"]\& C(X)\,.
\end{tikzcd}}
Donc pour montrer que $i_{C(X)}=C(i_X)$ dans $\mathrm{Cat}_{/S}W_S^{d\, -1}$, il suffit de prouver que $r$ est dans $W_S^d$. Mais on sait que le foncteur $C$ \uline{transforme} $W_S$ en $W_S^d$, donc $C(i_X)$ est dans $W_S^d$, et \inserted{$r\circ C(i_X)=\mathrm{id}$,} \deleted{$\unsure{C(i_X)}$} $r$ est dans $W_S^d$, OK !
% archive sheet number: 123
\scanpage{125}\mspage{70}
J'ai envie de grouper ce qui a été prouvé dans un énoncé récapitulatif.

\uline{Théorème}\inserted{-\uline{Scholie}} \quad On considère le diagramme de catégories de la page 8.\textsuperscript{*}
% slanted
\margin{On suppose sur $W$ les conditions a) b) c) de p.\,\unsure{4 pr.|4 sq.|4}, plus L 5 (p.\,\unsure{43}). Cf. Commentaires p.\,\unsure{81}-\deleted{\unsure{82}}}

\textcircled{I} Quand on munit toutes les 13 catégories du diagramme du localiseur $W_S$ ($W$-équivalences ``grossières'', i.e. $W$-équivalence de $\mathrm{Cat}$, sans plus), \uline{tou\deleted{\unsure{tes|te}}s} les foncteurs du diagramme induisent des équivalences sur les catégories localisées.\textsuperscript{*}
% slanted
\margin{De plus, si $\unsure{\mathcal F|F}$ est une sous-catégorie pleine de $\mathrm{Cat}_{/S}$ qui contient une des deux sous-catégories $\text{\uline{Fib}}_0(\unsure{S})$, $\text{\uline{Cofib}}_0 S$, \deleted{$\unsure{\mathrm{Parf}\,S}$} (ce qui inclut le cas de 8 catégories du diagramme, autres que $\mathrm{Cat}_{/S}$ \deleted{\unsure{elle}}-même), le foncteur $(\mathcal F,W_S\mid\mathcal F)$ dans $(\mathrm{Cat}_{/S},W_S)$ induit une équiv. pour les localisées.}
Désignons par
\[
 \mathrm{HOT}^{\unsure{\mathrm{lc}|\ell\mathrm{c}}}_W(S)
\]
``la'' catégorie localisée commune.

\textcircled{II} Munissons les quatre catégories \deleted{\unsure{bordantes}} gauches
\[
 \mathrm{Cat}_{/S}\hookleftarrow\text{\uline{Liss}}\,S\hookleftarrow\text{\uline{Fib}}\,S\longleftarrow\text{\uline{Fibsc}}\,S
 \tag*{$(*)$}
\]
du localiseur $W_S^d$ de la $W$-équivalence fine \uline{droite} sur $S$ (
\diagram{\begin{tikzcd}[ampersand replacement=\&,column sep=tiny,row sep=tiny]
 f:X\arrow[rr]\arrow[dr,no head]\&{}\&Y\arrow[dl,no head]\\
 {}\&S\&{}
\end{tikzcd}}
est dans $W_S^d$ ssi les ${}_{s\backslash}X\unsure{\longrightarrow|\text{---}}{}_{s\backslash}Y$ le sont, pour tout $s\in\mathrm{Ob}\,S$). Les foncteurs du diagramme i.e. ceux de $(*)$ entre ces catégories, \deleted{\unsure{et}} \inserted{donc aussi} leurs composés, induisent des équivalences entre les localisées.\textsuperscript{*}
% slanted
\margin{Plus généralement, pour toute sous-catégorie pleine $\mathcal F$ de $\mathrm{Cat}_{/S}$ qui contient $\text{\uline{Fib}}\,S$ (et dualement, avec $W_S^g$ au lieu de $W_S^d$, si $\mathcal F$ contient $\text{\uline{Cofib}}\,S$).}
Désignons par
\[
 \mathrm{HOT}_W(S)
\]
``la'' catégorie localisée commune.

Enoncé dual pour
\[
 \mathrm{Cat}_{/S}\hookleftarrow\text{\uline{Prop}}\,S\hookleftarrow\text{\uline{Cofib}}\,S\hookleftarrow\text{\uline{Cofibsc}}\,S,
 \tag*{$(*')$}
\]
et les localiseurs $W_S^g$ dans ces catégories ($f\in W_S^g$ ssi les $f_{/s}:X_{/s}\longrightarrow Y_{/s}$ sont dans $W$). On trouve une ``même catégorie localisée, notée
\[
 \mathrm{HOT}_W^\circ(S),
\]
et on \deleted{\unsure{a}} \deleted{\unsure{un}} tautologiquement
% archive sheet number: 124
\scanpage{126}\mspage{71}
\[
 \mathrm{HOT}_W^\circ(S)\approx\bigl(\mathrm{HOT}_W(S^\circ)\bigr)^{\deleted{\unsure{\circ}}}.
\]

\textcircled{III} Sur les six catégories internes du \unsure{diagramme|diagrammes}, les trois \deleted{\unsure{espèces de}} \inserted{notions de} $W$-\deleted{\unsure{équivalence}}, en termes de $W_S,W_S^d,W_S^g$ coïncident, et coïncident avec les $W$-équivalences par fibres \inserted{$W_S^f$}. Sur les \deleted{\unsure{trois}} \inserted{trois} catégories de $(*)$ autres que $\mathrm{Cat}_{/S}$, les localiseurs $W_S^f,W_S^d$ coïncident\textup{;} sur les trois catégories de $(*')$ autres que $\mathrm{Cat}_{/S}$, les localiseurs $W_S^f,W_S^g$ coïncident. \deleted{\unsure{L}}Mais en général,\textsuperscript{*}
% slanted
\margin{Sans doute dans tous les cas où $S$ n'est pas équivalente à \deleted{la \unsure{cat. ponctuelle}} \inserted{une catégorie} discrète.}
sur ces six catégories, \deleted{\unsure{cette}} ces localiseurs $W_S^f$ ne sont pas égaux à $W_S$, plus précisément donnent des localisées non équivalentes, \deleted{\unsure{et}} elles sont donc munies chacune de \uline{deux} localiseurs $\mathrm{ess}^{\mathrm{t}}$ distincts, $W_S^f$ et $W_S$. \deleted{\unsure{La|Sur} $\mathrm{Cat}_{/S}$} catégorie $\mathrm{Cat}_{/S}$ est la seule catégorie du diagramme sur laquelle les localiseurs $W_S,W_S^d$ et $W_S^g$ sont tous les trois $\mathrm{ess}^{\mathrm{t}}$ distincts. C'est la seule qui permette d'obtenir simultanément les trois catégories localisées, \unsure{comme|nommément}
\diagram{\scriptsize\begin{tikzcd}[ampersand replacement=\&,column sep=tiny,row sep=large]
 {}\&\mathrm{HOT}^{\unsure{\mathrm{lc}|\ell\mathrm{c}}}_W(S)\approxeq(\mathrm{Cat}_{/S})W_S^{-1}
 \arrow[dl,leftrightarrow,"{\deleted{\unsure{i}}\unsure{\rho|p}}"']\arrow[dr,"{\unsure{\sigma|j}}"]\&{}\\
 \mathrm{HOT}_W(S)\approxeq\mathrm{Cat}_{/S}(W_S^d)^{-1}\&{}\&
 \mathrm{HOT}^\circ_W(S)\approxeq\deleted{\unsure{(}}(\mathrm{Cat}_{/S})(W_S^g)^{-1}\,.
\end{tikzcd}}
Ici, $\rho$ et $\sigma$ sont des foncteurs de localisation (comme il est évident sur la définition).

Tout foncteur ou foncteur composé du diagramme page 8, allant d'une catégorie interne au diagramme vers une catégorie \unsure{bordante}, induit un foncteur pour les localisées pour $W_S^d$ (quand la catégorie but fait partie des quatre gauches,
% archive sheet number: 125
\scanpage{127}\mspage{72}
rappelées dans $(*)$ p.~70), soit pour $W_S^g$ (quand ces catégories font partie des quatre du côté droit, rappelées dans $(*')$), lesquels \deleted{\unsure{équiva}} localiseurs coïncident encore avec $W_S^f$ sauf dans le seul cas de la catégorie $\mathrm{Cat}_{/S}$. Tous les foncteurs ainsi obtenus se réduisent, à équivalence près, à deux, suivant que la catégorie localisée but est $\mathrm{HOT}_W(S)$ (cas de \deleted{\unsure{foncteur} \illegible{1} du côté} \inserted{la localisation} $W_S^d$), ou qu'elle est $\mathrm{HOT}_W^\circ(S)$ (cas de la localisation $W_S^g$). On trouve ainsi des foncteurs (en sens inverse de $\unsure{\rho|p}$ et de $\sigma$)
\diagram{\begin{tikzcd}[ampersand replacement=\&,row sep=large]
 {}\&\mathrm{HOT}_W^{\unsure{\mathrm{lc}|\ell\mathrm{c}}}(S)\arrow[dl,"i"']\arrow[dr,"j"]\&{}\\
 \mathrm{HOT}_W(S)\&{}\&\mathrm{HOT}_W^\circ(S)\,.
\end{tikzcd}}
Les foncteurs $\deleted{\unsure{R}}\rho,\sigma$ s'interchangent par dualité $S\mapsto S^\circ$, et de même pour $i,j$.

\textcircled{IV} Les foncteurs $\rho,\deleted{\unsure{j}}i$ sont adjoints l'un de l'autre dans cet ordre, i.e. on a pour deux objets \deleted{$X,Y$ de $\mathrm{Cat}$} $\xi\in\mathrm{Ob}\,\mathrm{HOT}_W(S)$, $\unsure{\eta|y}\in\mathrm{Ob}\,\mathrm{HOT}^{\mathrm{lc}}_W(S)$
\[
 \mathrm{Hom}(\rho(\deleted{\unsure{\eta}}\xi),\deleted{\unsure{\xi}}\eta)\simeq
 \mathrm{Hom}(\deleted{\unsure{\eta}}\xi,\deleted{\unsure{j}}i(\deleted{\unsure{\xi}}\eta))\deleted{\unsure{,}}\quad\deleted{\unsure{\xi}}
\]
\deleted{d'où \unsure{il}} \deleted{\unsure{et}} (\unsure{isom|isom.} de bifoncteurs). Et dualement les foncteurs $\sigma,j$ \deleted{\unsure{,}} sont adjoints l'un de l'autre dans cet ordre\,: pour $\xi\in\mathrm{Ob}\,\mathrm{HOT}^\circ_W(S)$, $\eta\in\mathrm{Ob}\,\mathrm{HOT}^{\mathrm{lc}}_W(S)$, on a
\[
 \mathrm{Hom}(\xi,j(\eta))\simeq\mathrm{Hom}(\sigma(\xi),\eta).
\]
% archive sheet number: 126
\scanpage{128}\mspage{73}
Il s'ensuit notamment que \uline{les foncteurs} $i,j$ \uline{sont pleinement fidèles}. \deleted{Pour que l'\unsure{élém$^{\mathrm t}$} $\xi$ de $\mathrm{HOT}_W(S)$, défini par l'objet $X$ de $\mathrm{Cat}_{/S}$ soit dans l'image essentielle de $i$, il \unsure{f. et s.} \unsure{que}}

\textcircled{V} \deleted{\unsure{Des}}On a des foncteurs évidents
\[
 \mathrm{HOT}_W(S)\unsure{\longrightarrow|\text{---}}\text{\uline{Hom}}\,(S^\circ,\mathrm{Hot}_W)
 \tag*{$(*)$}
\]
\[
 \mathrm{HOT}^\circ_W(S)\longrightarrow\text{\uline{Hom}}\,(S,\mathrm{Hot}_W),
 \tag*{$(*')$}
\]
le premier s'interprète sur les \deleted{trois}\inserted{quatre} catégories $(*)$ \inserted{indifféremment} \deleted{$\text{\uline{Fibsc}}\,S$} (bordantes gauches) de \unsure{page|p.}~70, \deleted{\unsure{qui}} servant à définir $\mathrm{HOT}_W(S)$, par
\[
 s\longrightarrow\mathrm{hot}_W(X_s)
\]
dans le cas des trois catégories distinctes de $\mathrm{Cat}_{/S}$, et par
\[
 s\longmapsto\mathrm{hot}_W({}_{s\backslash}X)
\]
dans le cas de $\mathrm{Cat}_{/S}$. La dépendance fonctorielle par rapport à $s$ est évidente dans le cas de $\mathrm{Cat}_{/S}$ (donné par $s\mapsto\mathrm{hot}_W({}_{s\backslash}X)$), et dans les cas de $\text{\uline{Fibsc}}\,S$, $\text{\uline{Fib}}\,S$, où on dispose, pour $u:s\unsure{\text{---}|\longrightarrow}t$ dans $S$, du changement de base
\[
 u^*:X_{\unsure{b|t}}\text{---}X_s.
\]
Dans le cas de $\text{\uline{Liss}}\,S$, la dépendance fonctorielle se déduit de celle \deleted{\unsure{du}} de $\mathrm{hot}_W({}_{s\backslash}X)$, grâce au fait que l'inclusion
\[
 X_s\longrightarrow{}_{s\backslash}X
\]
est alors une $W$-équivalence.

On a une interprétation duale de l'application
% archive sheet number: 127
\scanpage{129}\mspage{74}
$(*')$, \unsure{en termes|en écrivant}
\[
 s\longmapsto\mathrm{hot}_W(X_s)\qquad\text{ou}\qquad s\longmapsto\mathrm{hot}({}_{s\backslash}X),
\]
avec la loi fonctorielle en $s$ déduite des foncteurs ${}_{t\backslash}X\unsure{\text{---}|\longrightarrow}{}_{s\backslash}X$ (pour $u:s\text{---}t$) dans le cas de $\mathrm{Cat}_{/S}$, des foncteurs \unsure{cochangement} de base dans les cas de $\text{\uline{Cofib}}\,\deleted{\unsure{S}}$ et $\text{\uline{Cofibsc}}\,S$, et de la variante pour $\mathrm{Cat}_{/S}$ et de la $W$-équivalence
\[
 X_s\longrightarrow X_{/s}
\]
dans le cas \deleted{\unsure{de}} $\text{\uline{Prop}}\,S$.

On trouve alors ceci :
% slanted
\margin{il faudra encore écrire la \unsure{vérification}}
Pour qu'un objet $\xi$ de $\mathrm{HOT}_W(S)$ (resp. $\mathrm{HOT}^\circ_W(S)$) soit dans l'image essentielle de $\mathrm{HOT}^{\unsure{\mathrm{lc}|\ell\mathrm{c}}}_W(S)$, il f. et s. que le foncteur associé $S^\circ\text{---}\mathrm{Hot}_W$ (resp. $S\text{---}\mathrm{Hot}_W$) soit un système local${}^*$, i.e. transforme \unsure{toutes|toute|les} flèches de $S$ en un iso. \inserted{i.e. se factorise par $(\Pi\,\unsure{S})^\circ\longrightarrow\mathrm{Hot}_W$ resp. $(\Pi S)^{\deleted{\unsure{\circ}}}\text{---}\mathrm{Hot}_W$}
% slanted
\margin{où $\Pi S$ est le groupoïde fondamental de $S$}
Donc pour \unsure{un} $X$ dans $\mathrm{Cat}_{/S}$, cela signifie soit que les foncteurs
\[
 {}_{\unsure{t}\backslash}X\longrightarrow{}_{s\backslash}X
\]
associés à des $u:s\unsure{\longrightarrow|\text{---}}t$ dans $S$ sont tous des $W$-équivalences (cas de $\mathrm{HOT}_W(S)$), soit que les foncteurs
\[
 X_{/s}\text{---}X_{/t}
\]
le soient (cas de $\mathrm{HOT}^\circ_W(S)$), (NB. \deleted{\unsure{les}} Ces conditions ne sont nullement équivalentes. Mais si $X$ lisse $S$,
% archive sheet number: 128
\scanpage{130}\mspage{75}
la première implique la seconde, et l'inverse est vrai si $X$ est propre sur $S$, et dans ce cas, \deleted{\unsure{les}}la condition équiva\deleted{\unsure{lent}}ut à ce que $X$ soit \deleted{\unsure{dans|des}} $\text{\uline{Fibhot}}\,S$. Et \unsure{enfin|en fait}, si $X$ est déjà dans $\text{\uline{Fib}}\ \text{\uline{hot}}\,S$, alors les deux conditions sont automatiquement satisfaites.)
% slanted
\margin{\unsure{On en est donc|On en vient donc} dans le diagramme \unsure{à} introduire aussi $\text{\uline{Liss}}_0 S,\ \text{\uline{Prop}}_0 S$ !}

% slanted
\margin{le diagramme des \uline{quinze} catégories \unsure{remarquables} !}
\diagram{\begin{tikzcd}[ampersand replacement=\&,row sep=2pt,column sep=1pt,font=\scriptsize]
 \text{\uline{Fibsc}}_0\,S\arrow[dd,hook]\arrow[dr]\&{}\&{}\&{}\&{}\&{}\&\text{\uline{Cofibsc}}_0\,S\arrow[dd,hook]\arrow[dl]\\
 {}\&\deleted{\unsure{F}}\text{\uline{Fib}}_0\,S\arrow[dd,hook]\arrow[dr,hook]\&{}\&\mathrm{Parf}\,S\arrow[dl,hook]\arrow[dd,hook]\arrow[dr,hook]\&{}\&\text{\uline{Cofib}}_0\,S\arrow[dl,hook]\arrow[dd,hook]\&{}\\
 \text{\uline{Fibsc}}\,S\arrow[dr,hook]\&{}\&\text{\uline{Liss}}_0\,S\arrow[dd,hook]\arrow[dr,hook]\&{}\&\text{\uline{Prop}}_0\,S\arrow[dd,hook]\arrow[dl,hook]\&{}\&\text{\uline{Cofibsc}}\,S\arrow[dl,hook]\\
 {}\&\text{\uline{Fib}}\,S\arrow[dr,hook]\&{}\&\text{\uline{Fibhot}}\,S\arrow[dd,hook]\arrow[uuulll,dashed,"\Phi_0" description]\arrow[uuurrr,dashed,"\Psi_0" description]\&{}\&\text{\uline{Cofib}}\,S\arrow[dl,hook]\&{}\\
 {}\&{}\&\text{\uline{Liss}}\,S\arrow[dr,hook]\&{}\&\text{\uline{Prop}}\,S\arrow[dl]\&{}\&{}\\
 {}\&{}\&{}\&\mathrm{Cat}_{/S}\arrow[uuulll,dashed,bend left=20,"\Phi"']\arrow[uuurrr,dashed,bend right=20,"\Psi"']\&{}\&{}\&{}.
\end{tikzcd}}

De plus, si deux éléments $\xi\in\mathrm{Ob}\,\mathrm{HOT}_W(S)$, $\xi'\in\mathrm{HOT}^\circ_W(S)$ proviennent d'un même $\unsure{y|\eta}\in\mathrm{Ob}\,\mathrm{HOT}^{\unsure{\mathrm{lc}|\ell\mathrm{c}}}_W(S)$, alors les deux foncteurs \deleted{\unsure{associés} \unsure{sur}} correspondants
\[
 (\Pi S)^\circ\xrightarrow{\varphi}\mathrm{Hot}_W\xleftarrow{\varphi'}\Pi S
\]
sont \uline{associés}, i.e. \deleted{chacun est} que chacun est isomorphe à ``l'inverse'' de l'autre.
\[
 \left\{\begin{aligned}
 \varphi^\circ(s)&=\varphi(s)\\
 \varphi^\circ(u)&=\varphi(u)^{-1},
 \end{aligned}\right.
\]

\textcircled{VI} (NB Dans la nouvelle version du diagramme, où j'ai rajouté des catégories $\text{\uline{Liss}}_0 S$ et $\text{\uline{Prop}}_0 S$ pour rendre le \unsure{truc|tracé} plus transparent, les catégories appelées ``bordantes'' deviennent le ``rez de chaussée'', et les catégories ``internes'', le ``premier étage'' ou ``étage''. Tous les \inserted{\unsure{7}} carrés du diagramme sont cartésiens. Toutes les flèches sont des inclusions pl. fid, sauf les \deleted{\unsure{carrés}} flèches horizontales \unsure{extérieures} des catégories de Cat-fibrés ou de Cat-cofibrés \uline{scindés}, avec morphismes exclusivement cartésiens ou cocartésiens et \unsure{mêmes|même} compatibles avec les scindages, vers les cat. corresp. \unsure{admettant|donnant} des structures non scindées)
% archive sheet number: 129
